\bookStart{The Icelandic Rune Poem}
\setBookCode{IcelandicRunePoem}

\begin{flushright}%
\textbf{Dating:} Mediæval.%TODO

\textbf{Meter:} Unclear.
\end{flushright}%

\section{Introduction}

The \textbf{Icelandic Rune Poem} (abbrev.~\IcelandicRunePoem) is one of two West Norse rune poems describing the Younger Futhark.

\subsection{Textual preservation and transmission}

\IcelandicRunePoem\ is preserved in several manuscripts, none mediæval.  The oldest is AM 687 d 4to (henceforth 687d), dating to ca.~1500 \ce, where the poem takes up the leaf 1v.  In this ms.~each rune takes up one line.  The names of the runes are not given, but the name is written out by the rune, so that the first stanza begins \emph{ᚠ er frę́nda róg}, the second \emph{ᚢ er skýgia grátur}, and so forth.  687d/1v is not physically fragmented, but much of the ink has faded and large parts of thep oem are therefore illegible.

TODO.  It is only attested in late manuscripts which often have major disagreements with each other.

\subsection{Meter and stanza}

The meter of \IcelandicRunePoem\ is a divergent form of \Ljodahattr.  In particular, the a-verse always contains three lifts (certainly owing to the obligatory rune-name), and the c-verse may end in a dip following a heavy lift.
The poem is strictly structured.  Each rune is given two lines, the first consisting of two half-verses (the a- and b-verse), the second being a form of c-verse.  These two lines will henceforth be referred to as stanzas.

\subsection{Style and contents}

Like in the two other rune poems, the 16 runes in \IcelandicRunePoem\ are described by reference to the object to which their name refers (\emph{fé} ‘wealth’ for ᚠ, et c.).  Like its Scandinavian sister-poem \textlink{NorwegianRunePoem}, \IcelandicRunePoem\ uses kennings to describe these objects.  Numerous such kennings are shared, sometimes word-for-word, between the two.
Like in the two other rune poems, the very first word in each stanza is the respective rune’s name, which is a noun.  The name is then followed by exactly three kennings, one in each verse.


TODO: reference Mikael Males article about rune poems


\sectionline

\section{Text}

\bvg\bva%
\alst{F}é es \edtrans{\alst{f}rę́nda róg}{strife of kinsmen}{\Bfootnote{Self-explanatory; wealth causes conflict between heirs.  The same expression is found in \textlink{NorwegianRunePoem}.}} \hld\ ok \edtrans{\alst{f}lǿðar viti}{beacon of the flood}{\Bfootnote{A formulaic type of \textsc{gold}-kenning (TODO: \parencite{Meissner1921}); since gold does not rust, it is the fire which shines when submerged in water.}} &
\ind ok \edtrans{\alst{g}raf-sęiðs \alst{g}ata}{grave-saithe’s \ken{serpent’s} street}{\Bfootnote{Another formulaic kenning.  The serpent or Wyrm that lies upon its hill of treasures is a prominent image in Old Germanic legend; cf.~particularly Fathomer (\textlink{Reginsmal}, \textlink{Fafnismal}, and \VolsungaSaga) and the unnamed dragon in \textlink{Beowulf}.}}.\eva

\bvb Wealth is strife of kinsmen and beacon of the flood \\
\ind and grave-saithe’s \ken{serpent’s} street.\evb\evg


\bvg\bva%
Úr es \alst{sk}ýja grátr \hld\ ok \alst{sk}ára þvęrrir &
\ind ok \alst{h}irðis \alst{h}atr.\eva

\bvb Drizzle is weeping of clouds and ... \\
\ind and shepherd’s hatred.\evb\evg


\bvg\bva%
Þurs es \alst{k}venna \alst{k}vǫl \hld\ ok \alst{k}letta í·búi &
\ind ok \alst{v}arð-rúnar \alst{v}err.\eva

\bvb Thurse is women’s torment and indweller of hills \\
\ind and husband of the cairn-whisperess \ken{trollwoman}.\evb\evg


\bvg\bva%
\alst{Ǫ̇}ss es \alst{a}ldinn gautr \hld\ ok \alst{Ǫ̇}s·garðs jǫfurr, &
\ind ok \alst{V}al-hallar \alst{v}ísi.\eva

\bvb Os is ancient Geat, and Osyard’s chief, \\
\ind and Walhall’s overseer.\evb\evg


\bvg\bva%
Ręið es \alst{s}itjandi \alst{s}ę́la \hld\ ok \alst{s}núðig fęrð &
\ind ok \alst{jó}s \alst{ę}rfiði.\eva

\bvb Chariot is sitting bliss and twirling journey \\
\ind and horse’s heavy work.\evb\evg


\bvg\bva%
Kaun es \alst{b}arna \alst{b}ǫl \hld\ ok \alst{b}ar-dagi &
\ind ok \alst{h}old-fúa \alst{h}ús.\eva

\bvb Boil is children’s curse and TODO \\
\ind and house of flesh-rot.\evb\evg


\bvg\bva%
Hagall es \alst{k}alda \alst{k}orn \hld\ ok \alst{k}nappa drífa &
\ind ok \alst{s}náka \alst{s}ótt.\eva

\bvb Hail is cold kernel and storm of beads \\
\ind and sickness of snakes.\evb\evg


\bvg\bva%
Nauð es \alst{þ}ýjar \alst{þ}rǫ́ \hld\ ok \alst{þ}ungr kostr &
\ind ok \alst{v}ás-samlig \alst{v}erk.\eva

\bvb Need is maidservant’s yearning and scant choice \\
\ind and working in wet-cold weather.\evb\evg


\bvg\bva%
\alst{Í}ss es \alst{á}ar bǫrkr \hld\ ok \alst{u}nnar þękja &
\ind ok \alst{f}ęigra manna \alst{f}ár.\eva

\bvb Ice is river’s bark and wave’s roof \\
\ind and fey men’s danger.\evb\evg


\bvg\bva%
Ár es \alst{g}umna \alst{g}óði \hld\ ok \alst{g}ótt sumar &
\ind \emph{ok} \alst{a}l-gróinn \alst{a}kr.\eva

\bvb Year is men’s boon and good summer \\
\ind (and) all-grown acre.\evb\evg


\bvg\bva%
Sól es \alst{sk}ýja \alst{sk}jǫldr \hld\ ok \alst{sk}ínandi rǫðull &
\ind ok \alst{í}sa \alst{a}ldr-tregi.\eva

\bvb Sun is the shield of clouds and shining wheel \\
\ind and ice-sheets’ life-sorrow.\evb\evg


\bvg\bva%
Týr es \alst{ęi}n-hęndr \alst{ǫ̇}ss \hld\ ok \alst{u}lfs lęifar &
\ind ok \alst{h}ofa \alst{h}ilmir.\eva

\bvb Tew is the one-handed Os and the wolf’s leftovers \\
\ind and lord of hoves.\evb\evg


\bvg\bva%
Bjarkan es \alst{l}aufgat \alst{l}im \hld\ ok \alst{l}ítit tré &
\ind ok \alst{u}ng-samligr \alst{v}iðr.\eva

\bvb Birch is leafy branch and little tree \\
\ind and youthful wood.\evb\evg


\bvg\bva%
\alst{M}aðr es \alst{m}anns gaman \hld\ ok \alst{m}oldar auki &
\ind ok \alst{sk}ipa \alst{sk}ręytir.\eva

\bvb Man is man’s joy and the product of dust \\
\ind and adorner of ships.\evb\evg


\bvg\bva%
Lǫgr es \alst{v}ellanda \alst{v}atn \hld\ ok \alst{v}íðr kętill &
\ind ok \alst{g}lǫmmungr \alst{g}rund.\eva

\bvb Liquid is boiling water and wide kettle \\
\ind and TODO.\evb\evg


\bvg\bva%
Ýr es \alst{b}ęndr bogi \hld\ ok \alst{b}rot-gjarnt járn &
\ind ok \alst{f}ęnju \alst{f}lęygir.\eva

\bvb Yew is a bent bow and easily broken iron \\
\ind and arrow’s hurler.\evb\evg

\sectionline
